```latex
\documentclass[11pt]{article}

\usepackage{amsmath,amssymb,amsthm,mathtools}
\usepackage[margin=1in]{geometry}
\usepackage{hyperref}

\newtheorem{theorem}{Theorem}
\newtheorem{lemma}{Lemma}
\newtheorem{proposition}{Proposition}
\newtheorem{remark}{Remark}

\newcommand{\Hf}{\mathbb{H}}
\newcommand{\boxplusop}{\boxplus}
\newcommand{\inv}[1]{#1^{-1}}

\title{A Cubic Consequence of a Ternary Hyperadditive Relation (in the Quotient Case)}
\author{}
\date{}

\begin{document}
\maketitle

\begin{abstract}
Let $\Hf$ be a hyperfield with multivalued addition $\boxplusop$. The question asks whether
$0\in x\boxplusop y\boxplusop z$ implies
$0\in x^3\boxplusop y^3\boxplusop z^3\boxplusop(-xyz)\boxplusop(-xyz)\boxplusop(-xyz)$
for arbitrary hyperfields.
We give a complete proof for the important class of \emph{Krasner quotient hyperfields} $K/T$
(i.e.\ hyperfields obtained from a field $K$ by quotienting by a multiplicative subgroup $T\le K^\times$),
for which the implication holds by lifting to $K$ and using the classical factorization
$x^3+y^3+z^3-3xyz=(x+y+z)(x^2+y^2+z^2-xy-yz-zx)$.
We also recall that not every hyperfield is a Krasner quotient hyperfield.
\end{abstract}

\section{Preliminaries}

We use the standard axiomatic notion of hyperfield as in the literature, e.g.\ \cite[Def.\ 2.1]{KKS2025}.
Given elements $a,b\in \Hf$, the hyperaddition $a\boxplusop b$ is a nonempty subset of $\Hf$.
We extend $\boxplusop$ to finitely many summands by associativity (which holds as an axiom of hyperfields).

A basic consequence of the axioms is the \emph{reversibility} law
\[
z\in x\boxplusop y \iff x\in z\boxplusop (-y),
\]
see e.g.\ \cite[Prop.\ 2.3(v),(vi)]{KKS2025}.

\section{Krasner quotient hyperfields}

We recall Krasner's quotient construction (see \cite{Krasner1983} and also \cite[\S2]{KKS2025}).
Let $K$ be a field and $T\le K^\times$ a multiplicative subgroup.
Write $[a]$ for the $T$--coset of $a\in K$, with the convention $[0]=0$.
The quotient set $K/T=\{0\}\cup (K^\times/T)$ becomes a hyperfield with multiplication
$[a]\cdot[b]=[ab]$ and hyperaddition
\begin{equation}\label{eq:quotadd}
[a]\boxplusop [b] \;:=\; \bigl\{\, [a+b t]\;\big|\; t\in T\,\bigr\},
\end{equation}
see \cite{Krasner1983,KKS2025}. We call such a hyperfield a \emph{Krasner quotient hyperfield}.

\begin{lemma}[Lifting a ternary zero-sum in a quotient hyperfield]\label{lem:lift}
Let $K$ be a field, $T\le K^\times$, and let $\Hf:=K/T$.
Let $x=[a],y=[b],z=[c]\in \Hf$ (with $a,b,c\in K$).
If $0\in x\boxplusop y\boxplusop z$ in $\Hf$, then there exist $t_1,t_2,t_3\in T$ such that
\[
a t_1 + b t_2 + c t_3 = 0
\qquad\text{in }K.
\]
In particular, there exist representatives $a'\in x$, $b'\in y$, $c'\in z$ with $a'+b'+c'=0$ in $K$.
\end{lemma}

\begin{proof}
By associativity, $0\in (x\boxplusop y)\boxplusop z$.
Thus there exists $u\in x\boxplusop y$ such that $0\in u\boxplusop z$.
In a hyperfield, $0\in u\boxplusop z$ implies $u=-z$ by reversibility and uniqueness of additive inverses;
equivalently $-z\in x\boxplusop y$.

Unwinding the definition \eqref{eq:quotadd}, the containment $-z\in x\boxplusop y$ means:
there exists $t_2\in T$ such that
\[
[a+b t_2]=[-c].
\]
Hence $a+b t_2\in (-c)T$, so there exists $t_3\in T$ such that
\[
a + b t_2 = -c t_3,
\]
i.e.\ $a + b t_2 + c t_3=0$.
Taking $t_1=1\in T$ gives the first claim. Setting $a'=a t_1$, $b'=b t_2$, $c'=c t_3$
gives representatives in the respective cosets with $a'+b'+c'=0$.
\end{proof}

\section{The classical cubic identity}

\begin{lemma}[Ring identity]\label{lem:ring}
Let $R$ be a commutative ring and let $a,b,c\in R$ satisfy $a+b+c=0$.
Then
\[
a^3+b^3+c^3-3abc = 0.
\]
\end{lemma}

\begin{proof}
The standard factorization holds in $\mathbb{Z}[X,Y,Z]$:
\[
X^3+Y^3+Z^3-3XYZ = (X+Y+Z)\bigl(X^2+Y^2+Z^2-XY-YZ-ZX\bigr).
\]
Substituting $X=a$, $Y=b$, $Z=c$ in the commutative ring $R$ yields
\[
a^3+b^3+c^3-3abc = (a+b+c)\bigl(a^2+b^2+c^2-ab-bc-ca\bigr)=0,
\]
since $a+b+c=0$.
\end{proof}

\section{Main theorem for quotient hyperfields}

\begin{theorem}\label{thm:quotientcase}
Let $\Hf$ be a Krasner quotient hyperfield $K/T$.
For any $x,y,z\in \Hf$, if $0\in x\boxplusop y\boxplusop z$, then
\[
0\in x^3\boxplusop y^3\boxplusop z^3\boxplusop (-xyz)\boxplusop (-xyz)\boxplusop (-xyz).
\]
\end{theorem}

\begin{proof}
Write $x=[a]$, $y=[b]$, $z=[c]$ with $a,b,c\in K$.
Assume $0\in x\boxplusop y\boxplusop z$.
By Lemma~\ref{lem:lift}, there exist representatives $a'\in x$, $b'\in y$, $c'\in z$
such that
\[
a'+b'+c'=0 \qquad\text{in }K.
\]
Applying Lemma~\ref{lem:ring} in the field $K$ gives
\[
a'^3 + b'^3 + c'^3 - 3a'b'c' = 0 \qquad\text{in }K.
\]
Now interpret this equality in the quotient hyperfield.
We have
\[
x^3 = [a'^3],\qquad y^3=[b'^3],\qquad z^3=[c'^3],
\]
and
\[
xyz = [a'b'c']\quad\text{so}\quad -xyz = [-a'b'c'].
\]
By the definition of hyperaddition in a quotient hyperfield, the hyper-sum
\[
[a'^3]\boxplusop [b'^3]\boxplusop [c'^3]\boxplusop [-a'b'c']\boxplusop[-a'b'c']\boxplusop[-a'b'c']
\]
contains the class of the ordinary sum
\(
a'^3+b'^3+c'^3+(-a'b'c')+(-a'b'c')+(-a'b'c'),
\)
simply by choosing the subgroup multipliers $t=1$ at each hyperaddition step.
But that ordinary sum is exactly $a'^3+b'^3+c'^3-3a'b'c'=0$ in $K$.
Hence the coset $[0]=0$ lies in the displayed hyper-sum.
Translating back to $x,y,z$ gives the desired containment.
\end{proof}

\section{Remarks on the general (non-quotient) case}

\begin{remark}
Theorem~\ref{thm:quotientcase} applies to all hyperfields that arise as $K/T$ for a field $K$.
This includes the main hyperfields used in tropical and real algebraic geometry (e.g.\ the sign and tropical hyperfields),
which are of quotient type.

However, not every hyperfield (in the axiomatic sense) is isomorphic to a Krasner quotient hyperfield of a field:
Massouros constructed non-quotient hyperfields in \cite{Massouros1985}.
Thus Theorem~\ref{thm:quotientcase} does not, by itself, settle the implication for arbitrary hyperfields.
\end{remark}

\bibliographystyle{plain}
\begin{thebibliography}{10}

\bibitem{Krasner1983}
M.~Krasner.
\newblock A class of hyperrings and hyperfields.
\newblock {\em International Journal of Mathematics and Mathematical Sciences}, 6
  (1983), Article ID 240850.

\bibitem{Massouros1985}
Ch.~G. Massouros.
\newblock Methods of constructing hyperfields.
\newblock {\em International Journal of Mathematics and Mathematical Sciences},
  8(4):725--728, 1985.

\bibitem{KKS2025}
D.~E. K{\k e}dzierski, K.~Kuhlmann, and H.~Stoja{\l}owska.
\newblock Algorithmic construction of real hyperfields from minimal axioms.
\newblock arXiv:2508.19418v1, 2025.

\end{thebibliography}

\end{document}
```